\section{Limitations}
The frozen analysis input retains one fake-labelled WELFake row with an exact lightly normalized body match in ISOT: both become empty under artifact stripping and escape the builder's stripped-hash rule. Thus the residual is not completely free of exact cross-corpus overlap. Both identifiers occur in the historical exclusion list used for duplicate-removal checks. File hashes and counts establish input identity and size, not independent overlap verification or article veracity. Inspection of all saved paired-validation selection manifests found neither identifier in any training or evaluation selection across the eight seeds and three models. This finding applies to this pair and does not establish that all residual overlap has been removed.

The principal analysis covers three English news corpora dominated by 2016--2017 US politics, one lexical model and two fine-tuned transformers. The WELFake residual is constructed rather than independently collected. Dataset labels were not re-verified by human fact-checkers, and may encode source reliability rather than article-level veracity. Celebrity adds a separately collected external target after lexical overlap screening, not an independent fourth training corpus; no temporal holdout was evaluated. The paired run reuses the original document universe and eight seed identifiers, and its protocol was developed after the archived findings were known; it is a methodological validation, not an independent-data or prospectively registered confirmation.

Pairing holds article identities fixed for masking comparisons within each length-matching setting. L versus C0 changes article selection intentionally and uses target labels. Training is subsampled, and the transformers see only 256 tokens. Deletion before truncation changes input content; matching whole-word deletion counts does not match subword counts or positions. Source expressions were developed after inspecting predictive cues, and entity masking uses a small NER model. Neither transformation removes all information about its intended cue family. None establishes exclusive causal attribution.

In the three-corpus validation, only TF-IDF received the full 13-condition paired design; the transformers received eight conditions. Archived source splits, temporal masking, source-only transformer contrasts and duplicate-removal checks were not repeated with document pairing. The original duplicate-removal rerun did not cover RoBERTa or the added conditions. Saved provenance checks do not independently reproduce raw-text matching, NER or GPU training. Archived and new runs also differ in software, so changes between their estimates cannot be assigned solely to the pairing fix. Some original environment records are incomplete.

All seeds sample from one fixed universe. Seed-level intervals describe variability under this protocol, rather than uncertainty across future datasets. The six-contrast Holm summaries are descriptive because the family was chosen after earlier exploration. Pair-level analyses are unadjusted and exploratory, and the original dataset-size rule was narrowly missed.

Celebrity is small and domain-specific, with author-supplied labels and provisional numeric story groups. The audit records shared-content overlap and an unresolved cross-corpus label conflict, and one item lacks a body; group-based exclusions are conservative but cannot eliminate all semantic overlap. Fact-checking language and long historical articles remain. Target-side length matching uses labels, and its 428 articles differ from the 496-article primary target. All models share these fixed targets across seeds, so seed variability does not quantify article/story sampling uncertainty. Per-article prediction scores were not exported, preventing independent metric recalculation or a paired story bootstrap from the saved outputs. The external protocol followed earlier findings and was not formally registered; primary/secondary status is local to this extension, not a new confirmatory claim.
